\section{堆的虚拟地址区域规划}

内核堆需要一段专属连续虚拟地址。在第~\ref{ch:high-half}~章高半区布局基础上细分：

\begin{figure}[H]
  \centering
  % figures/ch10/fig01-kernel-vaddr-layout.tex — auto-extracted from chapter source
\begin{tikzpicture}[
  block/.style={draw, minimum width=7cm, font=\small, anchor=south west},
  addr/.style={font=\ttfamily\scriptsize\color{gray}, anchor=east},
]
  \node[block, minimum height=1.6cm, fill=red!10] (dm) at (0, 0)
    {\parbox{6.5cm}{\centering\textbf{直接映射区}\\
     \texttt{virt = phys + 0xC0000000}}};
  \node[block, minimum height=1.2cm, fill=yellow!15] (alloc) at (0, 1.6)
    {\parbox{6.5cm}{\centering\textbf{VMM 动态分配区}\\
     \texttt{vmm\_alloc\_pages} bump}};
  \node[block, minimum height=1.8cm, fill=green!15] (heap) at (0, 2.8)
    {\parbox{6.5cm}{\centering\textbf{内核堆}\\
     \texttt{kmalloc}/\texttt{kfree}，初始 64\,KiB，按需 grow}};
  \node[block, minimum height=0.8cm, fill=gray!10] (unmapped) at (0, 4.6)
    {\parbox{6.5cm}{\centering 未映射 / 保留}};
  \node[addr] at (-0.2, 0)    {\hex{C0000000}};
  \node[addr] at (-0.2, 1.6)  {\hex{D0000000}};
  \node[addr] at (-0.2, 2.8)  {\hex{E0000000}};
  \node[addr] at (-0.2, 4.6)  {\hex{F0000000}};
  \node[addr] at (-0.2, 5.4)  {\hex{FFFFFFFF}};
  \draw[decorate, decoration={brace, amplitude=5pt, mirror}]
    (7.2, 2.8) -- (7.2, 4.6) node[midway, right=8pt, font=\scriptsize] {256\,MiB};
\end{tikzpicture}

  \caption{内核虚拟地址空间细分布局}
  \label{fig:kernel-vaddr-layout}
\end{figure}

\texttt{vmm.h} 中定义的三个宏：

\codefile{src/include/kernel/vmm.h}
\begin{lstlisting}[style=cstyle]
#define KHEAP_START     0xE0000000u
#define KHEAP_MAX_SIZE  (256u * 1024u * 1024u)    /* 256 MiB */
#define KHEAP_END       (KHEAP_START + KHEAP_MAX_SIZE)  /* 0xF0000000 */
\end{lstlisting}

\begin{designbox}
堆区域必须和 \texttt{vmm\_alloc\_pages} 的地址空间分开。后者内部用 bump allocator
（\texttt{g\_next\_vaddr} 单调递增），如果堆也用同一空间，两个分配器会互相踩踏。
隔离到不同区间后各管各的，互不干扰。
\end{designbox}

% ============================================================================

